\chapter{Ereditarietà e polimorfismo in Java}
\label{chap:ereditarieta-polimorfismo-java}

Dopo la realizzazione di classi e degli algoritmi su array, si
introducono l'ereditarietà, le conversioni tra riferimenti e il
polimorfismo. Questi meccanismi permettono di riutilizzare e
specializzare comportamenti già definiti
(\cref{chap:realizzazione-classi-java,chap:ricerca-ordinamento-java}).%
\footnote{Materiale di riferimento: S. Canazza,
\emph{Java: Ereditarietà},
\nolinkurl{2026_FI_14_canazza_Java_Ereditarieta.key.pdf}.}

\section{Ordinamento di oggetti}
\label{sec:fi-java-ereditarieta-ordinamento-oggetti}

Gli algoritmi di ordinamento della
\cref{chap:ricerca-ordinamento-java} si applicano anche a oggetti,
quando è disponibile un criterio di confronto. A differenza dei
numeri, l'ordine tra oggetti dipende dal tipo e dall'operazione
definita dalla classe (\cref{sec:fi-java-cicli-oggetti}).

Per esempio, \texttt{String.compareTo()} confronta due stringhe
secondo l'ordinamento lessicografico dei valori UTF-16. Restituisce
un intero negativo, nullo o positivo quando la prima stringa
precede, equivale o segue la seconda secondo tale criterio.
Per ordinare stringhe mediante \texttt{selectionSort()},
il confronto numerico \lstinline|v[i] < min| viene sostituito da
\lstinline|v[i].compareTo(min) < 0|. Il procedimento e il numero di confronti restano invariati;
per le stringhe, tuttavia, il costo di ciascun confronto
dipende dalla loro lunghezza.

Per altri oggetti, come \texttt{BankAccount}, occorre stabilire
quale proprietà usare per il confronto. L\'interfaccia
\texttt{Comparable} è approfondita nella
\cref{sec:fi-java-interfacce-comparable}. La possibilità di
riutilizzare comportamenti comuni fra classi conduce
all'ereditarietà.

\section{Ereditarietà tra classi}
\label{sec:fi-java-ereditarieta}

\begin{definition}[Ereditarietà]
\label{def:fi-java-ereditarieta}
L'\emph{ereditarietà} permette di definire una classe più
specializzata a partire da una classe esistente. La nuova classe,
detta \emph{sottoclasse}, riutilizza i comportamenti accessibili della
\emph{superclasse} e può aggiungere o ridefinire membri.
\end{definition}

\subsection{Superclasse e sottoclasse}
\label{sec:fi-java-ereditarieta-conti}

La classe \texttt{BankAccount} della
\cref{sec:fi-java-classi-incapsulamento} gestisce il saldo con
\texttt{deposit()}, \texttt{withdraw()} e \texttt{getBalance()}.
Un conto di risparmio usa le stesse operazioni e aggiunge
\texttt{addInterest()} per accreditare gli interessi.

\begin{example}[Conto di risparmio]
La classe \texttt{SavingsAccount} estende \texttt{BankAccount}
introducendo il tasso annuo e il metodo per accreditare gli
interessi:

\begin{unipdcode}[language=Java]
public class SavingsAccount extends BankAccount {
    private final double interestRate;

    public SavingsAccount(double rate) { interestRate = rate; }
    public void addInterest() { deposit(getBalance() * interestRate / 100); }
}
\end{unipdcode}


\texttt{addInterest()} usa \texttt{getBalance()} e \texttt{deposit()}
della superclasse. Il campo \texttt{balance} resta privato e non è
accessibile direttamente dalla sottoclasse. La classe aggiunta
non duplica l'implementazione del conto bancario.
\end{example}

La sottoclasse rappresenta un caso particolare della superclasse:
ogni \texttt{SavingsAccount} è un \texttt{BankAccount}, mentre
non vale il contrario. La relazione descrive i tipi degli oggetti,
non il numero dei loro membri.

\subsection{Sintassi e gerarchie di ereditarietà}
\label{sec:fi-java-ereditarieta-sintassi}

La dichiarazione
\lstinline|class Sottoclasse extends Superclasse| stabilisce
l'ereditarietà. Una classe Java può estendere direttamente una
sola classe; in assenza di \texttt{extends}, la superclasse
implicita è \texttt{java.lang.Object}, ad eccezione di
\texttt{Object} stessa.

\begin{definition}[Superclasse universale \texttt{Object}]
\label{def:fi-java-ereditarieta-object}
\texttt{java.lang.Object} è la radice della gerarchia delle classi
Java. Le classi ne ereditano, direttamente o indirettamente,
i metodi accessibili, tra cui \texttt{toString()} ed
\texttt{equals()}, salvo eventuali sovrascritture. Il loro comportamento è
approfondito nella \cref{sec:fi-java-interfacce-object}.
\end{definition}

La \cref{fig:fi-java-ereditarieta-gerarchia} rappresenta
l'ereditarietà diretta e indiretta. Nei diagrammi UML la freccia
con triangolo vuoto è diretta verso la superclasse.

\begin{figure}[htbp]
  \centering
  \begin{tikzpicture}[unipd/diagram,x=1cm,y=1cm]
  \tikzset{fi/inheritance/.style={unipd/node,align=center,
    minimum width=3.1cm,minimum height=0.68cm,
    inner xsep=5pt,inner ysep=4pt,font=\small}}
  \node[fi/inheritance] (obj) at (0,2.6) {\texttt{Object}};
  \node[fi/inheritance] (bank) at (0,1.3) {\texttt{BankAccount}};
  \node[fi/inheritance] (save) at (0,0) {\texttt{SavingsAccount}};
  \draw[thick] (bank.north) -- (0,2.14);
  \draw[thick,fill=white] (-0.14,2.14) -- (0,2.26) -- (0.14,2.14) -- cycle;
  \draw[thick] (save.north) -- (0,0.84);
  \draw[thick,fill=white] (-0.14,0.84) -- (0,0.96) -- (0.14,0.84) -- cycle;
\end{tikzpicture}

  \caption{Gerarchia di ereditarietà delle classi
  \texttt{Object}, \texttt{BankAccount} e \texttt{SavingsAccount}.}
  \label{fig:fi-java-ereditarieta-gerarchia}
\end{figure}

\subsection{Metodi nuovi, ereditati e sovrascritti}
\label{sec:fi-java-ereditarieta-metodi}

Una sottoclasse può introdurre nuovi metodi, ereditare quelli
della superclasse oppure \emph{sovrascriverli}
(\emph{overriding}). In \texttt{SavingsAccount},
\texttt{addInterest()} è nuovo e \texttt{getBalance()} è
ereditato.

\begin{definition}[Sovrascrittura di un metodo]
\label{def:fi-java-ereditarieta-override}
Si ha \emph{sovrascrittura} quando una sottoclasse fornisce una
nuova implementazione di un metodo di istanza ereditato, con
firma compatibile e tipo di ritorno ammesso dalle regole di Java.
L'annotazione \texttt{@Override} permette al compilatore di
verificare che il metodo sovrascriva effettivamente un metodo
ereditato.
\end{definition}

Per applicare una commissione fissa ai versamenti,
\texttt{SavingsAccount} può sovrascrivere \texttt{deposit()}.
I membri seguenti si aggiungono alla classe dell'esempio
precedente:

\begin{unipdcode}[language=Java]
@Override
public void deposit(double amount) {
    super.deposit(amount);
    super.withdraw(FEE);
}

private static final double FEE = 2.58;
\end{unipdcode}


La parola \texttt{super} permette di invocare l'implementazione
del metodo nella superclasse: \texttt{super.deposit(amount)}
accredita l'importo senza richiamare il metodo sovrascritto.
La successiva \texttt{super.withdraw(FEE)} addebita la commissione.
Una chiamata non qualificata a \texttt{deposit(amount)}
all'interno di quel metodo richiamerebbe invece se stessa,
provocando una ricorsione senza caso base.

\texttt{super.nomeMetodo(argomenti)} non costruisce un altro
oggetto: seleziona il metodo della superclasse per l'istanza
corrente. Un metodo ereditato può comunque essere invocato
normalmente attraverso un riferimento alla sottoclasse.

Nei metodi di istanza, \texttt{this} indica l'oggetto corrente;
\texttt{super} permette di richiamare i membri della superclasse
accessibili alla sottoclasse. Entrambi operano sulla stessa
istanza. I metodi ereditati possono accedere ai campi
\texttt{private} dichiarati nella propria classe, mentre i metodi
della sottoclasse non possono accedervi direttamente
(\cref{sec:fi-java-classi-incapsulamento}).

\section{Conversioni tra riferimenti}
\label{sec:fi-java-ereditarieta-conversioni}

Un riferimento a una sottoclasse può essere assegnato a una
variabile del tipo di una sua superclasse (\emph{upcasting}).
La conversione è implicita e non modifica né copia l'oggetto:

\begin{unipdcode}[language=Java]
SavingsAccount risparmio = new SavingsAccount(10);
BankAccount conto = risparmio;
Object oggetto = risparmio;
\end{unipdcode}

Le tre variabili contengono riferimenti alla stessa istanza,
come mostra la \cref{fig:fi-java-ereditarieta-riferimenti}.
Il tipo dichiarato del riferimento limita i membri accessibili
in compilazione: \texttt{conto.deposit(100)} è valido, mentre
\texttt{conto.addInterest()} non lo è, perché
\texttt{addInterest()} non è dichiarato in \texttt{BankAccount}.

\begin{figure}[htbp]
  \centering
  \begin{tikzpicture}[unipd/diagram,x=1cm,y=1cm]
  \tikzset{fi/ref/.style={unipd/node,align=center,
    minimum width=2.4cm,minimum height=0.6cm,
    inner xsep=4pt,inner ysep=4pt,font=\scriptsize}}
  \node[fi/ref] (a) at (-3.0,1.15) {\texttt{risparmio}};
  \node[fi/ref] (b) at (-3.0,0.15) {\texttt{conto}};
  \node[fi/ref] (c) at (-3.0,-0.85) {\texttt{oggetto}};
  \node[unipd/node,align=center,minimum width=3.1cm,
    minimum height=1.25cm,font=\scriptsize]
    (inst) at (1.1,0.15) {Istanza di \texttt{SavingsAccount}\\
      \texttt{balance = 0}\\\texttt{interestRate = 10}};
  \draw[unipd/directed edge] (a.east) -- (inst.north west);
  \draw[unipd/directed edge] (b.east) -- (inst.west);
  \draw[unipd/directed edge] (c.east) -- (inst.south west);
\end{tikzpicture}

  \caption{Tre riferimenti di tipo diverso alla medesima
  istanza di \texttt{SavingsAccount}.}
  \label{fig:fi-java-ereditarieta-riferimenti}
\end{figure}

\begin{example}[Sostituibilità nei parametri]
Un metodo con parametro di tipo \texttt{BankAccount} può
ricevere un oggetto di qualunque sua sottoclasse:

\begin{unipdcode}[language=Java]
public static void accredita(BankAccount conto, double importo) {
    conto.deposit(importo);
}

public static void esempio() {
    SavingsAccount risparmio = new SavingsAccount(10);
    accredita(risparmio, 100);
}
\end{unipdcode}

Il parametro \texttt{conto} riceve per valore una copia del
riferimento, non un nuovo oggetto
(\cref{sec:fi-java-file-parametri}). Se
\texttt{deposit()} è sovrascritto, viene eseguita la versione
della classe effettiva dell'oggetto.
\end{example}

La conversione inversa (\emph{downcasting}) richiede un cast
esplicito ed è valida soltanto se l'oggetto è un'istanza del tipo
richiesto o di una sua sottoclasse:

\begin{unipdcode}[language=Java]
SavingsAccount altro = (SavingsAccount) conto;
altro.addInterest();
\end{unipdcode}

Il cast non modifica il tipo effettivo dell'oggetto. Se il
riferimento indica un oggetto incompatibile, l'esecuzione
provoca \texttt{ClassCastException}, un'eccezione non controllata
(\cref{sec:fi-java-file-checked}).


\section{Polimorfismo}
\label{sec:fi-java-ereditarieta-polimorfismo}

\begin{definition}[Polimorfismo per sottotipi]
\label{def:fi-java-ereditarieta-polimorfismo}
Il \emph{polimorfismo per sottotipi} permette di utilizzare
oggetti di classi derivate attraverso riferimenti al tipo della
superclasse. Per un metodo di istanza sovrascritto, l'implementazione
eseguita dipende dal tipo effettivo dell'oggetto.
\end{definition}

\begin{example}[Selezione dinamica dei metodi]
Si considerino \texttt{BankAccount} e \texttt{SavingsAccount},
che sovrascrive \texttt{deposit()} applicando una commissione
di \texttt{2.58} a ogni versamento.

\begin{unipdcode}[language=Java]
BankAccount[] conti = {
    new BankAccount(),
    new SavingsAccount(10)
};

for (BankAccount conto : conti) {
    conto.deposit(100);
    System.out.println(conto.getBalance());
}
\end{unipdcode}

\begin{unipdterminal}
100.0
97.42
\end{unipdterminal}

In entrambe le iterazioni, \texttt{conto} ha tipo dichiarato
\texttt{BankAccount}, ma si riferisce a oggetti di classi
differenti. La stessa invocazione \texttt{deposit(100)}
esegue quindi due implementazioni diverse: quella della
superclasse nel primo caso e quella sovrascritta nella
sottoclasse nel secondo.

La selezione avviene durante l'esecuzione, senza richiedere
conversioni esplicite o controlli sul tipo dell'oggetto.
\end{example}

Il \emph{tipo statico} del riferimento determina quali metodi
possono essere invocati; il \emph{tipo dinamico} dell'oggetto
determina quale implementazione di un metodo sovrascritto viene
eseguita. Questo meccanismo è detto \emph{late binding}
(selezione dinamica). Si distingue dall'\emph{early binding}
dei metodi sovraccarichi
(\cref{sec:fi-java-classi-overloading}), per i quali il compilatore
sceglie la firma applicabile in base ai tipi statici degli argomenti.

La selezione dinamica descritta riguarda i metodi di istanza
sovrascrivibili; non si applica allo stesso modo ai metodi
\texttt{static}, ai campi o ai costruttori.

\subsection{Verifica del tipo con \texttt{instanceof}}
\label{sec:fi-java-ereditarieta-instanceof}

L'operatore \texttt{instanceof} restituisce \texttt{true} se
l'oggetto indicato dal riferimento è compatibile con il tipo
specificato. Il risultato dipende dal tipo effettivo dell'oggetto,
non dal solo tipo dichiarato del riferimento; per \texttt{null}
è sempre \texttt{false}.

\begin{unipdcode}[language=Java]
BankAccount conto = new SavingsAccount(10);
if (conto instanceof SavingsAccount) {
    SavingsAccount risparmio = (SavingsAccount) conto;
    risparmio.addInterest();
}
\end{unipdcode}

Il controllo protegge il cast verso \texttt{SavingsAccount}:
quando la condizione è vera, la conversione non provoca
\texttt{ClassCastException}. Evita inoltre di assumere che ogni
\texttt{BankAccount} sia un conto di risparmio.
